\documentclass[11pt,a4paper]{article}
\usepackage[T1]{fontenc}
\usepackage[utf8]{inputenc}
\usepackage{lmodern,amsmath,amssymb}
\usepackage[margin=25mm]{geometry}
\usepackage{longtable,booktabs,array,calc}
\usepackage[hidelinks]{hyperref}
\hypersetup{pdftitle={The mass gap as a finite list of theorems: the Hamiltonian lattice route},pdfauthor={navstokgap research working notes}}
\newcommand{\tightlist}{\setlength{\itemsep}{0pt}\setlength{\parskip}{0pt}}
\setlength{\emergencystretch}{3em}
\setcounter{secnumdepth}{0}
\begin{document}
\hypertarget{the-mass-gap-as-a-finite-list-of-theorems-the-hamiltonian-lattice-route}{%
\section{The mass gap as a finite list of theorems: the Hamiltonian
lattice
route}\label{the-mass-gap-as-a-finite-list-of-theorems-the-hamiltonian-lattice-route}}

On the Hamiltonian lattice route the Jaffe--Witten conjecture decomposes
into four named statements, and their present status is: (T1) every
finite periodic lattice has a positive gap \(\Delta_{a,L}\) above a
unique gauge-invariant ground state, proved below with all constants
explicit; (T2) at strong coupling the gap stays bounded below uniformly
in the volume, established in the Euclidean form by the
Osterwalder--Seiler cluster expansion and cited here with its
hypotheses, while the naive Kato bound fails because the magnetic term
is extensive; (T2\('\)) the gap of the infinite lattice stays positive
for every coupling, which is the statement that four-dimensional
non-abelian lattice gauge theory has no Coulomb phase, open, and false
for the abelian group by Guth and Fröhlich--Spencer; (T3) along the
asymptotic-freedom curve \(g(a)\) the infinite-volume gap \(\Delta_a\)
satisfies \(\Delta_a/(\hbar c\Lambda(g(a)))\to m/\Lambda\in(0,\infty)\),
open, and this single limit is the Millennium content once the theory
exists; (T4) the continuum, infinite-volume theory with the
Osterwalder--Schrader axioms, known in finite volume with an ultraviolet
cutoff removed (Balaban; Magnen--Rivasseau--Sénéor) and open beyond. The
small-volume expansion controls one corner: for \(L\ll\hbar c/\Lambda\)
the gap is \(\delta_1g(L)^{2/3}\hbar c/L\) (C133, Lüscher), so the
dimensionless gap \(z(L)=\Delta(L)L/(\hbar c)\) must cross over from
\(\propto g^{2/3}\) at small \(L\) to linear growth \(mL/(\hbar c)\) at
large \(L\). T2 in Hamiltonian form, a volume-uniform gap for
\(g\ge g_0(N)\), is closed in \href{strong-coupling-uniform-gap.md}{the
strong-coupling note}; the radius \(g_0\) is the number the entire
problem asks to push to zero. Sources are cited inline with reading
labels; nothing here is promoted.

\hypertarget{the-finite-lattice-hamiltonian}{%
\subsection{1. The finite-lattice
Hamiltonian}\label{the-finite-lattice-hamiltonian}}

Let \(G\) be a compact connected simple Lie group with Lie algebra basis
\(T^a\) normalized by \(\operatorname{tr}T^aT^b=\tfrac12\delta^{ab}\) in
a fixed faithful representation (for \(SU(N)\) the fundamental). Let
\(\Lambda=(a\mathbb Z/Na\mathbb Z)^3\) be the periodic spatial lattice
with spacing \(a\), side \(L=Na\), link set \(\mathcal E\) and plaquette
set \(\mathcal P\); \(|\mathcal E|=3N^3\), \(|\mathcal P|=3N^3\). The
Hilbert space is
\(\mathcal H=L^2(G^{\mathcal E},\mathrm{Haar}^{\otimes\mathcal E})\). On
the copy of \(G\) attached to a link \(\ell\), the electric field
components \(E^a_\ell\) are \(-i\) times the left-invariant vector
fields, and \(\sum_a(E^a_\ell)^2=-\Delta_\ell\) is the Laplace--Beltrami
operator of the bi-invariant metric, whose eigenvalues on the
Peter--Weyl blocks are the quadratic Casimirs \(C_2(R)\) with
multiplicity \((\dim R)^2\). The Kogut--Susskind Hamiltonian (Phys.
Rev.~D 11 (1975) 395, metadata) in the convention used here is
\[H=\frac{\hbar c}{a}\Big[\frac{g^2}{2}\sum_{\ell\in\mathcal E}(-\Delta_\ell)
+\frac{2}{g^2}\sum_{p\in\mathcal P}\big(N-\operatorname{Re}\operatorname{tr}U_p\big)\Big]
=\frac{\hbar c}{a}\,\big[g^2K+g^{-2}V\big],\] with \(g\) dimensionless,
\(U_p\) the ordered product of the four link variables around \(p\), and
\(K\), \(V\) the dimensionless electric and magnetic operators. Other
conventions rescale \(g^2\) by a fixed number. The dimensionless
spectrum of \(aH/(\hbar c)\) depends on \(g\), \(N\) and \(G\) only:
\[\Delta_{a,L}=\frac{\hbar c}{a}\,\delta(g;N,G).\]

Gauge transformations \(\gamma:\Lambda\to G\) act by
\((\mathcal V_\gamma\psi)(U)=\psi(\gamma_xU_\ell\gamma_y^{-1})\) for
\(\ell=(x,y)\); they are unitary, commute with \(H\) (the Laplacian is
bi-invariant and \(\operatorname{tr}U_p\) is conjugation invariant) and
preserve positivity of functions. The physical space
\(\mathcal H_{\rm phys}\) is the joint fixed space (Gauss's law).

\hypertarget{t1-every-finite-lattice-has-a-gap-above-a-unique-physical-ground-state}{%
\subsection{2. T1: every finite lattice has a gap above a unique
physical ground
state}\label{t1-every-finite-lattice-has-a-gap-above-a-unique-physical-ground-state}}

\textbf{Theorem 1.} For every \(a>0\), \(N\ge2\), \(g>0\) and \(G\) as
above:

\begin{enumerate}
\def\labelenumi{\arabic{enumi}.}
\tightlist
\item
  \(H\) is self-adjoint, bounded below, and has compact resolvent; its
  spectrum is a sequence \(E_0<E_1<\cdots\to\infty\) of eigenvalues of
  finite multiplicity.
\item
  \(E_0\) is simple with a strictly positive eigenfunction \(\psi_0\),
  so \(\Delta_{a,L}=E_1-E_0>0\).
\item
  \(\psi_0\in\mathcal H_{\rm phys}\), \(\mathcal H_{\rm phys}\) reduces
  \(H\), and the physical gap satisfies
  \(\Delta^{\rm phys}_{a,L}\ge\Delta_{a,L}>0\).
\end{enumerate}

\emph{Proof.} (1) \(-\Delta_\ell\) is essentially self-adjoint on
\(C^\infty(G)\), nonnegative, with compact resolvent, since \(G\) is a
compact manifold and Peter--Weyl gives the eigenvalues \(C_2(R)\) with
finite multiplicities and \(C_2(R)\to\infty\) along any enumeration of
irreducible representations. \(K=\sum_\ell(-\Delta_\ell)\) acts on the
finite tensor product with eigenvalues \(\sum_\ell C_2(R_\ell)\); for
each bound \(E\) only finitely many assignments \(\ell\mapsto R_\ell\)
have \(\sum_\ell C_2(R_\ell)\le E\), so \(K\) has compact resolvent.
\(V\) is multiplication by a continuous function with
\(0\le V\le2N|\mathcal P|\), hence bounded, and a bounded perturbation
of an operator with compact resolvent has compact resolvent by the
resolvent identity
\((H-z)^{-1}=(H_0-z)^{-1}[1+(g^{-2}V)(H_0-z)^{-1}]^{-1}\) for \(z\) far
in the negative axis. Discreteness with finite multiplicities and
\(E_n\to\infty\) follow, and \(E_0\) is an eigenvalue.

\begin{enumerate}
\def\labelenumi{(\arabic{enumi})}
\setcounter{enumi}{1}
\item
  The heat kernel of \(-\Delta\) on a compact connected Riemannian
  manifold is strictly positive for \(t>0\) (standard, from the strong
  maximum principle or the Peter--Weyl expansion; metadata level), so
  the kernel of \(e^{-tK}\) on \(G^{\mathcal E}\), a product of strictly
  positive kernels, is strictly positive. With Brownian motion \(X_t\)
  on \(G^{\mathcal E}\) generated by \(K\) (time rescaled by
  \(g^2\hbar c/a\)), the Feynman--Kac formula gives
  \[(e^{-tH}\psi)(U)=\mathbb E^U\Big[\exp\Big(-\frac{\hbar c}{ag^2}\int_0^tV(X_s)\,ds\Big)\psi(X_t)\Big],\]
  valid because \(V\) is bounded and continuous. For \(\psi\ge0\),
  \(\psi\ne0\), the weight is at least
  \(e^{-t\hbar c\cdot2N|\mathcal P|/(ag^2)}>0\) and the law of \(X_t\)
  has a strictly positive density, so \((e^{-tH}\psi)(U)>0\) for every
  \(U\): \(e^{-tH}\) is positivity improving. Reed--Simon IV, Theorem
  XIII.44 (metadata; number corroborated as in B78) then gives that
  \(E_0\) is simple with \(\psi_0>0\) almost everywhere, and \(E_1>E_0\)
  because \(E_1\) is the next eigenvalue of a discrete spectrum.
\item
  \(\mathcal V_\gamma\psi_0\) is a normalized eigenfunction for \(E_0\),
  hence \(\mathcal V_\gamma\psi_0=c(\gamma)\psi_0\) with \(|c|=1\); both
  functions are positive, so \(c(\gamma)=1\) and
  \(\psi_0\in\mathcal H_{\rm phys}\). The projection
  \(P=\int\prod_x d\gamma_x\,\mathcal V_\gamma\) onto
  \(\mathcal H_{\rm phys}\) commutes with \((H-z)^{-1}\), so
  \(\mathcal H_{\rm phys}\) reduces \(H\) and the restricted operator
  has compact resolvent, ground state \(\psi_0\), and next eigenvalue at
  least \(E_1\). \(\square\)
\end{enumerate}

Theorem 1 defines the object whose limits the conjecture is about. It
proves nothing about those limits: \(\Delta_{a,L}\) is \(\hbar c/a\)
times a number that may tend to zero as \(N\to\infty\) or as \(g\to0\).

\hypertarget{t2-strong-coupling-and-why-volume-uniformity-is-the-whole-difficulty}{%
\subsection{3. T2: strong coupling, and why volume uniformity is the
whole
difficulty}\label{t2-strong-coupling-and-why-volume-uniformity-is-the-whole-difficulty}}

At \(g\to\infty\) the operator \(g^2K\) dominates. Its ground state is
the constant function, with eigenvalue \(0\), and its excitations are
link representations. Gauge invariance excludes a single excited link:
the lowest physical excitation of \(K\) is a closed flux loop, at
minimum the four links of one plaquette in the lowest nontrivial
representation, with electric energy
\(4\cdot\tfrac12C_2(R_{\min})\,g^2\hbar c/a=2C_2(R_{\min})g^2\hbar c/a\);
for \(SU(N)\), \(C_2(\text{fund})=(N^2-1)/(2N)\). So the unperturbed
physical gap is
\[\gamma_{\rm phys}(g)=2C_2(R_{\min})\,g^2\,\frac{\hbar c}{a},\] uniform
in \(N\) because a single plaquette loop costs the same on every
lattice.

\textbf{The naive bound fails.} The magnetic term satisfies
\(0\le g^{-2}V\le2N|\mathcal P|g^{-2}\) in operator norm, so Kato's
perturbation bound gives
\(\Delta^{\rm phys}\ge\gamma_{\rm phys}-2\|g^{-2}V\|\hbar c/a =(2C_2g^2-4N|\mathcal P|g^{-2})\hbar c/a\),
positive only when \(g^4>2N|\mathcal P|/C_2\). The bound is extensive in
the volume and useless as \(N\to\infty\). Volume uniformity therefore
requires using that \(V\) is a sum of \(|\mathcal P|\) terms each of
norm at most \(2N g^{-2}\) and each acting on four links: locality, not
smallness of the norm.

\textbf{What is established.} Osterwalder and Seiler (Ann. Phys. 110
(1978) 440; abstract, B78) prove for the Euclidean lattice theory with
any compact gauge group, at sufficiently strong coupling, reflection
positivity, the existence of the infinite-volume limit and Wilson's
confinement bound by a cluster expansion; exponential clustering of
gauge-invariant correlations uniform in the volume, that is a mass gap
of the transfer matrix in units of \(1/a\), is the standard consequence
of the same expansion (not read at section level). Hamiltonian versions
rest on gap-stability theorems for local perturbations of commuting,
frustration-free Hamiltonians (Yarotsky, Commun. Math. Phys. 261 (2005)
799; Bravyi--Hastings--Michalakis, J. Math. Phys. 51 (2010) 093512; both
metadata): \(g^2K\) is a sum of commuting link terms with a product
ground state and a uniform local gap, and \(g^{-2}V\) is a sum of
bounded local terms. Their hypotheses assume finite-dimensional local
spaces or bounded local terms, and \(-\Delta_\ell\) is unbounded on the
infinite-dimensional \(L^2(G)\), so a Hamiltonian proof of T2 must
either truncate each link to the representations with
\(C_2(R)\le C_{\max}\) and control the truncation, or use the relative
boundedness of \(V\) with respect to \(K\). This is the first bounded
theorem in the list:

\begin{quote}
\textbf{T2 (Hamiltonian form, target).} There exist \(g_0(N)<\infty\)
and \(c(N)>0\) such that for all \(g\ge g_0\) and all lattice sizes
\(N\ge2\), \(\Delta^{\rm phys}_{a,L}\ge c(N)\,g^2\,\hbar c/a\).
\end{quote}

\hypertarget{t2-no-coulomb-phase-at-fixed-cutoff}{%
\subsection{\texorpdfstring{4. T2\('\): no Coulomb phase at fixed
cutoff}{4. T2\textquotesingle: no Coulomb phase at fixed cutoff}}\label{t2-no-coulomb-phase-at-fixed-cutoff}}

Fix \(a\) and \(g\) and let \(N\to\infty\). Define
\(\delta_\infty(g)=\liminf_{N\to\infty}\delta(g;N,G)\). T2 says
\(\delta_\infty(g)\ge c\,g^2>0\) for \(g\ge g_0\). The lattice form of
the conjecture at fixed cutoff is

\begin{quote}
\textbf{T2\('\).} \(\delta_\infty(g)>0\) for every \(g>0\).
\end{quote}

For \(G=U(1)\) this is false: Guth (Phys. Rev.~D 21 (1980) 2291) and
Fröhlich--Spencer (Commun. Math. Phys. 83 (1982) 411) prove a massless
Coulomb phase at weak coupling (abstract, B78), so
\(\delta_\infty(g)=0\) below a critical coupling. T2\('\) for a
non-abelian \(G\) is the statement that the transition the abelian
theory undergoes does not occur, that is confinement at all couplings,
which lattice computations support and no proof reaches. The heuristic
of \href{low-dimensional-mass-gap.md}{G07} §3 locates the difference:
the abelian theory is all valley, while the non-abelian potential has
flat directions of measure zero along which a transverse zero-point
energy grows. T2\('\) is where that mechanism would have to be turned
into an estimate at weak coupling.

\hypertarget{t3-the-scaling-limit-as-one-asymptotic-statement}{%
\subsection{5. T3: the scaling limit, as one asymptotic
statement}\label{t3-the-scaling-limit-as-one-asymptotic-statement}}

Along the continuum curve the bare coupling runs with the cutoff. With
\(b_0=11N/(48\pi^2)\) and \(b_1=34N^2/(3(16\pi^2)^2)\) for \(SU(N)\)
(the two-loop coefficients in the convention
\(\beta(g)=-b_0g^3-b_1g^5+\cdots\); Gross--Wilczek, Politzer, abstract,
B78), the lattice \(\Lambda\) parameter is defined by
\[a\Lambda_{\rm lat}(g)=(b_0g^2)^{-b_1/(2b_0^2)}\,e^{-1/(2b_0g^2)}\,[1+O(g^2)],\]
a dimensionless function that vanishes faster than any power of \(g\).

\begin{quote}
\textbf{T3.} The limit
\(\displaystyle\frac{m}{\hbar c\,\Lambda_{\rm lat}} =\lim_{g\to0}\frac{\delta_\infty(g)}{a\Lambda_{\rm lat}(g)}\)
exists and lies in \((0,\infty)\).
\end{quote}

This is the Jaffe--Witten mass gap \(m\) expressed on the lattice route:
the dimensionless infinite-volume gap must vanish as \(g\to0\) exactly
like \(e^{-1/(2b_0g^2)}\) times the prescribed power, with a positive
finite prefactor. Any weaker vanishing gives \(m=\infty\) in the
continuum (the theory is gapped only at the cutoff scale, as the compact
abelian theory in three dimensions is, G07 §4); any faster vanishing, or
\(\delta_\infty=0\) at some \(g\), gives \(m=0\). T3 presupposes
T2\('\). The order of limits is \(N\to\infty\) at fixed \(g\) first,
then \(g\to0\) with \(a=\Lambda_{\rm lat}^{-1}\,a\Lambda_{\rm lat}(g)\);
Jaffe and Witten note (p.~6, passage) that a gap uniform in the volume
would also serve to construct the infinite-volume limit, which is the
reverse order.

\textbf{The small-volume corner.} For \(L\ll\hbar c/\Lambda\) the
constant-mode sector of the continuum torus theory has gap
\(\delta_1g(L)^{2/3}\hbar c/L\) (C133, G07 Proposition 7), and Lüscher's
expansion (Nucl. Phys. B219 (1983) 233, abstract) supplies the
corrections in powers of \(g^{2/3}\) with \(g\) the renormalized
coupling at scale \(L\). In the variable \(z(L)=\Delta(L)L/(\hbar c)\)
the two ends read
\[z(L)=\delta_1\,g(L)^{2/3}\,[1+O(g^{2/3})]\ \ (L\to0),\qquad
z(L)\sim\frac{mL}{\hbar c}\ \ (L\to\infty),\] and Lüscher--Münster
(Nucl. Phys. B232 (1984) 445, abstract) place the crossover near
\(z\simeq2\). The conjecture, in this variable, is that \(z(L)\) is a
function that starts at zero and grows without bound, with \(m\) the
slope of its linear asymptote. The corner is the only place where the
gap of the continuum theory is presently computable, and it is
controlled by the mechanism of
\href{action-floor-yang-mills-gap.md}{G08}.

\hypertarget{t4-existence}{%
\subsection{6. T4: existence}\label{t4-existence}}

The lattice route needs the continuum limit of the Euclidean lattice
theory with Osterwalder--Schrader reconstruction. Known: ultraviolet
stability of four-dimensional lattice Yang--Mills in a finite volume
with the cutoff removed (Balaban, Commun. Math. Phys. 122 (1989) 355 and
the series it belongs to; Magnen--Rivasseau--Sénéor, Commun. Math. Phys.
155 (1993) 325, with an infrared cutoff; both metadata). Open: the
infinite-volume limit, the axioms in infinite volume, and any spectral
information. Jaffe and Witten (p.~7, passage) write that a construction
on a compact space alone would be a major step and that no present ideas
give a gap uniform in the volume.

\hypertarget{the-list}{%
\subsection{7. The list}\label{the-list}}

\begin{longtable}[]{@{}
  >{\raggedright\arraybackslash}p{(\columnwidth - 4\tabcolsep) * \real{0.3333}}
  >{\raggedright\arraybackslash}p{(\columnwidth - 4\tabcolsep) * \real{0.3333}}
  >{\raggedright\arraybackslash}p{(\columnwidth - 4\tabcolsep) * \real{0.3333}}@{}}
\toprule\noalign{}
\begin{minipage}[b]{\linewidth}\raggedright
Statement
\end{minipage} & \begin{minipage}[b]{\linewidth}\raggedright
Content
\end{minipage} & \begin{minipage}[b]{\linewidth}\raggedright
Status
\end{minipage} \\
\midrule\noalign{}
\endhead
\bottomrule\noalign{}
\endlastfoot
T1 & finite lattice: unique physical ground state, gap
\(\Delta_{a,L}=(\hbar c/a)\delta(g;N,G)>0\) & proved (Theorem 1) \\
T2 & \(\delta(g;N,G)\ge c(N)g^2\) for \(g\ge g_0(N)\), uniformly in
\(N\) & established: Euclidean form (Osterwalder--Seiler); Hamiltonian
form as a corollary of Yarotsky's theorem in
\href{strong-coupling-uniform-gap.md}{the strong-coupling note} \\
T2\('\) & \(\delta_\infty(g)>0\) for all \(g>0\) (no Coulomb phase) &
open; false for \(U(1)\) \\
T3 &
\(\delta_\infty(g)/(a\Lambda_{\rm lat}(g))\to m/(\hbar c\Lambda_{\rm lat})\in(0,\infty)\)
& open; the Millennium content given T4 \\
T4 & continuum infinite-volume theory with the axioms & finite-volume UV
stability known; rest open \\
S & small-volume corner \(z(L)=\delta_1g(L)^{2/3}[1+O(g^{2/3})]\) & C133
plus Lüscher's expansion \\
\end{longtable}

Three of the six rows are theorems or established results; the
conjecture is T2\('\) together with T3, given T4. The list replaces
``prove the mass gap'' by ``prove T2 in Hamiltonian form, then push
\(g_0\) down''.

\hypertarget{consequence-for-state}{%
\subsection{8. Consequence for STATE}\label{consequence-for-state}}

T2 in Hamiltonian form is now closed in
\href{strong-coupling-uniform-gap.md}{the strong-coupling note} as a
corollary of Yarotsky's gap-stability theorem, with threshold
\(g_0(N)=(48N^2/[(N^2-1)\beta_*])^{1/4}\); every later step is an
improvement of that threshold toward zero along the scaling curve. The
\(d=3\) target of STATE keeps its place as the case where T3 has no
transmutation to perform.
\end{document}
